\documentclass[10pt,a4paper]{article}
\usepackage[a4paper,margin=1.8cm]{geometry}
\usepackage[T1]{fontenc}
\usepackage{lmodern}
\usepackage{microtype}
\usepackage{booktabs}
\usepackage{tabularx}
\usepackage{array}
\usepackage[table]{xcolor}
\usepackage{enumitem}
\usepackage{titlesec}
\usepackage[colorlinks=true,linkcolor=blue!50!black,urlcolor=blue!50!black]{hyperref}

\setlist{nosep,leftmargin=1.3em}
\titlespacing*{\section}{0pt}{1.1ex}{0.6ex}
\titlespacing*{\subsection}{0pt}{0.9ex}{0.4ex}
\setlength{\parindent}{0pt}
\setlength{\parskip}{0.45ex}
\renewcommand{\arraystretch}{1.08}
\newcommand{\iss}[1]{\href{https://github.com/JieGH/RVV_code_gen_via_MLIR_xDSL/issues/#1}{\##1}}
\newcommand{\up}[1]{\textcolor{red!70!black}{#1}}
\newcommand{\down}[1]{\textcolor{blue!60!black}{#1}}

\title{\vspace{-1.2cm}\textbf{Counting the X60's L2 and prefetcher per rep:\\ what B3's slow state is}\\[2pt]
\large Brief of issue \iss{42} (implemented), with the first answer for \iss{41}}
\author{RVV\_code\_gen\_via\_MLIR\_xDSL, \texttt{pure-xdsl} branch, commit \texttt{5d04deb}}
\date{27 September 2026}

\begin{document}
\maketitle
\vspace{-0.8cm}

\paragraph{Scope.} \iss{41} asks why B3 (64$\times$384$\times$384) falls into a state about half as fast after a
short fast start (\iss{36}). Its first check found that the generic L2 and prefetch counters open on the X60 but
always read 0. \iss{42} is the spec that fixes this: find the X60's own event codes, count them per rep, and put
fast reps beside slow reps. This brief covers what was built, how it was checked, the review, and what the first
nine board launches show. All runs on \texttt{jbpi} (BananaPi, SpacemiT X60, VLEN=256), tile 16$\times$15,
\texttt{nc}=384 (the model's), \texttt{--no-nc-overrides}. \textbf{Every count here comes from the instrumented build}
(\texttt{INSTRUMENTED=1}, warm-up protocol, 1000 fixed reps): these numbers show where time goes, and are
\textbf{not speed numbers to quote} against the reference. Every log cited is committed (Appendix~A).

\section{Bottom line}
\begin{enumerate}
  \item \textbf{The X60's L2 and prefetch events can be counted.} They are raw codes \texttt{0xa0}--\texttt{0xbf} in
  the vendor kernel's perf tables; the board lets any code \texttt{0x00}--\texttt{0xff} through. All the ones we use
  count sensibly, checked on the board by a test that now runs with the suite.
  \item \textbf{New tooling:} \texttt{--raw-events x60-l2} (instrumented build only) adds a pack and a compute count
  per event to every rep in \texttt{reps\_*.txt}; \texttt{reps\_report.py} puts fast reps beside slow reps for every
  counter, over several launches; each launch records which core was busy during its timed run.
  \item \textbf{What B3's slow state is:} in slow reps the L2 is accessed \emph{exactly as often}, but demand L2
  misses are \up{2--4$\times$} and L1D prefetch hits and refills \down{about halve}. The L1D hardware prefetcher is
  doing about half its work, and the lines it no longer brings in become demand misses. The same pattern holds for
  both packers, for B3 with K=256, and for Reuse24's few slow reps.
  \item \textbf{It is not the shared L2 of cores 4--7, nor a neighbour.} On core 2 (the other cluster, its own L2) B3
  has two speeds just the same, with the same counter pattern; during a timed run no other core is busy.
  \item \textbf{Still open:} what makes the prefetcher back off, and what can stand in for it (next steps in
  \S5; the X60 has no prefetch hints, so the fixes specified are \iss{43} and \iss{44}).
\end{enumerate}

\section{What was built (\iss{42})}
\begin{tabularx}{\linewidth}{@{}>{\raggedright\arraybackslash}p{4.1cm}X@{}}
\toprule
Part & What it does \\
\midrule
Counter header & An optional second counter group of raw events (up to 8) beside the usual 10-event group, read
at the same pack and compute windows. Kept apart, so an event the board refuses never takes the usual counters down,
and a window the raw group was shared out for only spoils that rep's raw counts. Holds the checked
\texttt{x60-l2} preset. \\
Driver & \texttt{--raw-events name=code,...} or \texttt{x60-l2}: a \texttt{Pack\_<event>} and a
\texttt{Compute\_<event>} column per event at the end of the per-rep log; \texttt{nan} when the group cannot open
or a rep's window was shared out. Refused in the uninstrumented build, without \texttt{--save-reps}, and for a bad
list (not \texttt{name=code}, a bad name or code, a name twice, more than 8). The run overview says what was asked
for and whether it opened. \\
\texttt{compile\_macro\_bench.sh} & The same flag, refused before any board time without \texttt{--instr} or with
\texttt{--no-save-reps}. \texttt{provenance.txt} gains the raw events and each core's busy share \emph{during each
timed run} (read from \texttt{/proc/stat} on the board just before and just after \texttt{make run}). \\
\texttt{reps\_report.py} & A second table: per case and counter, the mean per rep over fast reps and over slow
reps of all launches given, and slow/fast. Same fast/slow split as the launch table (rep 0 left out), same
``fewer than 3 launches'' warning; the per-run line names the raw events and the busy share. \\
Tests & 4 counter-header tests, 13 driver tests, 6 report tests, 7 script tests, and an on-board check of the
preset (both groups open together, neither is shared out, each event counts as below). Full suite: 266 passed. \\
\bottomrule
\end{tabularx}

\section{The X60 events, checked on the board}
Counts per 64-byte line, streaming a 64\,KiB buffer (fits the 512\,KiB L2) and a 4\,MiB one (does not), 20
passes, core 4 (kernel 6.6.63 \texttt{\#2.2.7.2}; the firmware version is not readable without root).

\begin{center}\small
\begin{tabular}{@{}llrrrr@{}}
\toprule
event & code & 64\,KiB read & 4\,MiB read & 64\,KiB write & 4\,MiB write \\
\midrule
\texttt{l2\_load\_access} & 0xb8 & 1.03 & 1.01 & 0.86 & 1.01 \\
\texttt{l2\_load\_miss} & 0xb9 & 0.002 & 0.002 & 0.001 & 0.05 \\
\texttt{l2\_store\_access} & 0xbc & 1.02 & 0.99 & 0.86 & 1.01 \\
\texttt{l2\_ar\_channel\_stall\_cycle} & 0xbb & 0.00 & 2.84 & 0.03 & 28.7 \\
\texttt{l2\_aw\_channel\_stall\_cycle} & 0xbd & 0 & 0 & 0 & 0.001 \\
\texttt{l1d\_miss} & 0xab & 0.06 & 0.02 & 0.84 & 1.00 \\
\texttt{l1d\_prefetch\_refill} & 0xae & 0.93 & 0.94 & 0 & 0 \\
\texttt{l1d\_prefetch\_hit} & 0xaf & 1.93 & 1.96 & 0 & 0 \\
\bottomrule
\end{tabular}
\end{center}
One surprise: on a streaming \emph{read}, demand L2 misses stay near 0 even at 4\,MiB, because the prefetcher runs
ahead of the loads. The time spent waiting past the L2 shows up in the AR-channel stall cycles instead (0 at 64\,KiB,
2.8 per line at 4\,MiB). So a missing prefetch, not a bigger working set, is what turns into demand misses.

\section{Fast reps against slow reps}
Nine launches in the order V4 S4 V2 \textbar{} V2 S4 V4 \textbar{} V4 S4 V2 (V/S = vector/scalar packer, 4/2 =
core), so each pairing runs A-B-B-A-A-B, three launches each. Ratios are slow reps over fast reps, mean per rep,
pooled over the three launches; ``pack / compute'' are the two counter windows.

\begin{center}\small
\begin{tabular}{@{}lrcccc@{}}
\toprule
case (packer, core) & slow reps & cycles & L2 accesses & L2 load misses & L1D prefetch hits \\
 & & pack / compute & pack / compute & pack / compute & pack / compute \\
\midrule
B3 (vector, 4) & 51\% & 2.09 / 1.51 & 1.01 / 0.95 & \up{2.52 / 3.13} & \down{0.40 / 0.51} \\
B3 (vector, 2) & 64\% & 1.89 / 1.47 & 1.01 / 0.95 & \up{2.25 / 3.02} & \down{0.50 / 0.56} \\
B3\_K256 (vector, 4) & 60\% & 1.97 / 1.50 & 1.05 / 0.95 & \up{3.22 / 3.73} & \down{0.52 / 0.52} \\
B3\_K256 (scalar, 4) & 36\% & 2.10 / 1.46 & 1.27 / 0.93 & \up{4.21 / 3.39} & \down{0.31 / 0.49} \\
Reuse24 (vector, 4) & 7\% & 1.44 / 1.11 & -- & \up{1.62 / 1.23} & \down{0.62 / 0.48} \\
B3 (scalar, 4) & 0\% & \multicolumn{4}{c}{one speed in these three launches} \\
B2 (both packers, 4) & 0\% & \multicolumn{4}{c}{one speed} \\
\bottomrule
\end{tabular}
\end{center}

Also in B3's slow reps (vector, core 4): L1D prefetch refills 0.46 / 0.58, L1D misses in the compute window 3.06,
AR-channel stall cycles 0.43 / 0.48 (fewer requests go past the L2, while the core waits longer on the ones it makes).

\paragraph{How to read it.} The slow reps do the same work against the L2 (same access counts), so the L2 is not
filling up or thrashing: B3\_K256, whose B fits the L2, slows down more, not less. What changes is who brings the
lines in. In a fast rep the prefetcher does; in a slow rep it does about half as much, and each line it skips costs a
demand miss. Core 2 gives the same picture, so this is not the cores 4--7 cluster or anything running on cores 5--7.

\section{What is known, and what is next}
\begin{tabularx}{\linewidth}{@{}>{\raggedright\arraybackslash}p{5.3cm}X@{}}
\toprule
Ruled out & Evidence \\
\midrule
CPU clock, huge pages, B not fitting the L2 & \iss{36} \\
The L2 filling up or thrashing & same L2 access count in fast and slow reps; B3\_K256 slows more \\
The shared L2 of cores 4--7, a neighbour process & core 2 shows the same; busy share during a timed run: core 4
100\%, all others 0\% \\
\bottomrule
\end{tabularx}

\textbf{Named so far:} the L1D hardware prefetcher backs off to about half its work. \textbf{Next} (\iss{41}):
(1) the same counts in 100-rep windows around the switch, to see what changes just before it; (2) \texttt{ldb}
padding off 384 (\iss{41} step 4, \iss{38}); (3) stand-ins for the hardware prefetcher: the X60 has \emph{no} prefetch hints (no Zicbop, so
\texttt{prefetch.r} would do nothing), but it has \texttt{cbo.zero}, which runs in user mode on 64-byte lines. Two
fixes are now specified, on hold: \iss{43}, \texttt{cbo.zero} on the pack buffers so their writes stop reading
lines from the L2, and \iss{44}, packing B one column at a time (one stream instead of 15); (4) repeat on
\texttt{jbpi2}. Meanwhile B3's \texttt{nc}=128 override is back as the baseline.

\section{Review}
Two-axis review of the change (sub-agents, on standards and on the spec), then fixed before commit:
\begin{itemize}
  \item \emph{Standards:} no breach of a documented rule. One judgement call: the fast/slow split was written twice
  in \texttt{reps\_report.py}. \textbf{Fixed}: one function, used by both tables.
  \item \emph{Spec:} (a) the busy share was read over the whole launch, so it counted the board build (cores 1, 2, 6
  at 15--58\% in the nine launches). \textbf{Fixed}: read around each timed run; checked on a tenth launch (core 4
  at 100\%, every other core at 0\%). (b) Firmware version not recorded: \textbf{noted} as not readable without root.
  (c) The event table lacked per-row source and ``checked'' columns: \textbf{added} in \texttt{docs/metrics.md},
  with the three codes checked but left out of the preset. (d) The table in \iss{41}: \textbf{added} (\S8 there).
\end{itemize}
Caveat carried forward: the nine launches' busy shares predate fix (a), so they cannot show a neighbour during the
timed run; the core-2 result and the tenth launch are what rule one out.

\appendix
\section{Logs (all committed, \texttt{tests/pure\_xdsl\_blis\_rvv/logs/})}
Build \texttt{INSTRUMENTED=1}, warm-up protocol, 1000 reps, \texttt{--no-nc-overrides}, \texttt{--raw-events x60-l2},
\texttt{jbpi}, config \texttt{config\_issue41\_l2\_counters.json}.
\begin{itemize}
  \item core 4, vector: \texttt{run\_2026-09-27\_131302}, \texttt{run\_2026-09-27\_132606}, \texttt{run\_2026-09-27\_132843}
  \item core 4, scalar: \texttt{run\_2026-09-27\_131537}, \texttt{run\_2026-09-27\_132328}, \texttt{run\_2026-09-27\_133117}
  \item core 2, vector: \texttt{run\_2026-09-27\_131815}, \texttt{run\_2026-09-27\_132052}, \texttt{run\_2026-09-27\_133354}
  \item busy-share check (not in the tables): \texttt{run\_2026-09-27\_133825}
\end{itemize}
Beside them, in \texttt{.scratch/issues/}:
\begin{itemize}
  \item reports: \texttt{0025-reps-report-core4.md}, \texttt{0025-reps-report-core2.md}
  \item event probe: \texttt{0026-raw-probe.c}, \texttt{0026-raw-probe.txt}
  \item event table: \texttt{docs/metrics.md}, ``Raw events''
\end{itemize}

\end{document}
